% concatenated from SvininFunctionn2026.tex
\documentclass[reqno,a4paper,11pt]{amsart}
%%%%%%%%%%%%%%%%%%%%%%%%%%%%%%%%%%%%%%%%%%%%%%%%%%

\tolerance 2000
\textwidth 15cm
\textheight 23cm
\topmargin -.0cm
\oddsidemargin 0.5cm
\evensidemargin 0.5cm

\usepackage{amsmath,amstext,amsfonts,amsbsy,eucal,amssymb}
\usepackage[latin1]{inputenc}


\vfuzz2pt % Don't report over-full v-boxes if over-edge is small
\hfuzz2pt % Don't report over-full h-boxes if over-edge is small
\parskip 1mm


\numberwithin{equation}{section}

\newtheorem{theorem}{Theorem}[section]
\newtheorem{lemma}[theorem]{Lemma}
\newtheorem{sublemma}[theorem]{Sublemma}
\newtheorem{proposition}[theorem]{Proposition}
\newtheorem{corollary}[theorem]{Corollary}
\newtheorem{cor}[theorem]{Corollary}
\newtheorem{conjecture}[theorem]{Conjecture}
\newtheorem{fact}[theorem]{Fact}
\newtheorem{question}[theorem]{Question}

\theoremstyle{definition}
\newtheorem{definition}[theorem]{Definition}
\newtheorem{example}[theorem]{Example}
\newtheorem{remark}[theorem]{Remark}





\begin{document}

\parskip 4pt
\baselineskip 16pt

%%%%%%%%%%%%%%%%%%%%%%%%%%%%%%%%%
%%%%%%%%%%%%%%%%%%%%%%%%%%%%%%%%%

\title[On a certain arithmetic function defined via Bernoulli numbers]
{On a certain arithmetic function defined via Bernoulli numbers}

\author[Andrei K. Svinin]{Andrei K. Svinin}
\address{Andrei K. Svinin, 
Matrosov Institute for System Dynamics and Control Theory of 
Siberian Branch of Russian Academy of Sciences,
P.O. Box 292, 664033 Irkutsk, Russia}
\email{svinin@icc.ru}




\thanks{This work was carried out within the framework of  the state assignment of the Ministry of Education
and Science of the Russian Federation on the project No.  126021217175-3.
}



\date{\today}
%\subjclass[2000]{53D17, 37J35, 14H70, 14H40}

\keywords{Bernoulli number, Carmichael number, Giuga number}


% ----------------------------------------------------------------

\begin{abstract}
We investigate the integrality property of an arithmetic function defined on the set of odd numbers $n\geq 3$. It is constructed via Bernoulli numbers. As a result, we show that this function unifies three distinct number classes -- primes, Carmichael numbers, and Giuga numbers -- into a single integrality criterion. The article is elementary and therefore accessible to a broad audience.
\end{abstract}


\maketitle

\section*{Introduction}

The Bernoulli numbers $\left\{B_j\right\}$  are defined by the generating power series
\[
\frac{t}{e^t-1}=\sum_{j\geq 0} B_j\frac{t^j}{j !}.
\]

Next, we will use two theorems. The first is the classical von Staudt-Clausen theorem \cite{Clausen}, \cite{Hardy}, \cite{Staudt}.
\begin{theorem}
For any positive odd integer $n$:
\[
B_{n-1}+\sum_{p}\frac{1}{p}\;\;\mbox{is}\;\; \mbox{integer},
\]
where the summation runs over all primes  satisfying the condition $(p-1)\mid (n-1)$. 
\end{theorem}
This theorem completely describes the structure of the denominator of any Bernoulli number. 

Using Bernoulli numbers, we now define the function
\[
\chi(n)=\frac{\mathrm{denom}\left(\displaystyle\frac{B_{n-1}}{n-1}\right)}{n\cdot\mathrm{denom}\left(B_{n-1}+\displaystyle\frac{1}{n}\right)},
\]
for any odd integer $n\geq 3$.  Also we shall consider  $\tilde{\chi}(n)=n\cdot \chi(n)$. 

Our main goal in this paper is to investigate   integrality properties of $\chi(n)$.  It should be noted that the proofs of the statements given in the next section are quite elementary.

In the sequel, we shall  need the definitions of two types of composite numbers, namely Carmichael numbers and Giuga numbers. 
\begin{definition}
A composite positive integer $n$ is called a Carmichael number, if 
\[
a^{n-1}\equiv 1\left(\mathrm{mod}\; n\right)
\]
is valid for all $(a, n) = 1$.
\end{definition}
Informally speaking, Carmichael numbers are pseudoprimes that satisfy Fermat's little theorem for any base. 

It is, of course, difficult to verify directly from this definition whether a given number $n$ is a Carmichael number, but Korselt's theorem \cite{Crandall}, \cite{Korselt} then comes to our aid.
\begin{theorem}  (Korselt's criterion) 
A composite number $n$ is a Carmichael number if and only if it satisfies two conditions: $n$ is squarefree, and for every prime divisor $p$ of $n$:  $\left(p-1\right)\mid  \left(n-1\right)$. 
\end{theorem}  
In other words, we shall actually take this as the definition of a Carmichael number. We shall also need the definition of a Giuga number.
\begin{definition}
A Giuga number is a composite integer $n$ such that each prime divisor $p$ of $n$ satisfies 
\begin{equation}
p\mid \left(\frac{n}{p}-1\right) \Leftrightarrow p^2\mid (n-p) \Leftrightarrow n\equiv p\; (\mathrm{mod}\; p^2).
\label{867434}
\end{equation}
\end{definition}
Even Giuga numbers are well known and appear in the OEIS under the identifier \textrm{A007850}. This sequence begins $30, 858, 1722, 66198,\ldots$ However, we are only interested in odd Giuga numbers. It can be proven that any odd Giuga number (if it exists) is a Carmichael number, i.e., it satisfies the conditions of Korselt's criterion. 


If an odd Giuga number exists, it must be a product of at least 14 primes. The sharpest known lower bound to date,  implies that any odd Giuga number must have at least 19,908 decimal digits \cite{Borwein}.

The question of the existence of odd Giuga numbers is an open problem, closely related to the Agoh-Giuga conjecture. 
However, we do not wish to delve into this topic, as it is not the aim of this paper and refer the reader to authoritative sources, such as \cite{Borwein1}, \cite{Kellner1}.

%It is worth noting the paper \cite{Kellner2} by Kellner and Sondow, in which the author investigate the connections between Bernoulli numbers and Carmichael numbers. There they put forward the conjecture that the relation
%\[
%n\cdot \mathrm{denom}\left(n B_{n-1}\right)=\mathrm{denom}\left(B_{n-1}\right)
%\]
%holds if and only if $n$ is prime or a Carmichael number.


In what follows, we agree that $n$ denotes an odd integer greater than three, while $p$ denotes a prime number. We use  $B_{n-1}=A/D$ to denote  corresponding Bernoulli number in reduced form. Note that $A$ here  inherits the sign of $B_{n-1}$.

\section{Integrality properties  of the function $\chi(n)$}

In this section, we present several results concerning the integrality property  of the function $\chi(n)$ for different odd integers.
\begin{lemma}
For any prime $p\geq 3$, the function $\chi(p)$ is integer-valued.
\end{lemma}
\begin{proof}
By the von Staudt-Clausen theorem:
\[
\mathrm{denom}\left(B_{p-1}+\frac{1}{p}\right)=\frac{D}{p}.
\]
In addition, we take into consideration that
\[
\mathrm{denom}\left(\displaystyle\frac{B_{p-1}}{p-1}\right)=\frac{(p-1)D}{\gcd(p-1, \mathrm{num}(B_{p-1}) )}.
\]
Thus, 
\begin{equation}
\chi(p)=\frac{p-1}{\gcd(p-1, \mathrm{num}(B_{p-1}) )}.
\label{675498}
\end{equation}
The latter is clearly an integer.
\end{proof}
\begin{lemma}  \label{875643}
The value $\chi(p)$  is even for any prime $p$.
\end{lemma}
\begin{proof}
By the von Staudt-Clausen theorem, the denominator of $B_{p-1}$ is always divisible by 2. Hence, the numerator of $B_{p-1}$ is odd, while $p-1$ is evidently even. Therefore, $\chi(p)$ defined by     (\ref{675498}) is even for any prime $p$.
\end{proof}
\begin{lemma}
Let $n$ be a composite number that is not a Carmichael number. Then $\chi(n)$ is not an integer.
\end{lemma}
\begin{proof}
For any odd composite integer $n$, let $q$ be a prime divisor of $n$. Suppose that $q-1\nmid  n-1$.  Then $q\nmid  \mathrm{denom}\left(B_{n-1}/(n-1)\right)$, but $q\mid \mathrm{denom}\left(B_{n-1}+1/n\right)$. This means that in this case $\tilde{\chi}(n)=n\cdot \chi(n)$ is not an integer, and a fortiori neither is $\chi(n)$.   Thus, if we have found such a prime $q$, that is sufficient to prove that $\chi(n)$  is not an integer. 
\end{proof}
\begin{lemma}
For  any Carmichael number, the value of the function $\tilde{\chi}(n)$ is integer.  
\end{lemma}
\begin{proof}
To begin with, let us see how the expression for the function $\tilde{\chi}(n)$ can be written most conveniently, provided that its argument is a Carmichael number. 

Since $\gcd(A, D)=1$, then 
\[
\mathrm{denom}\left(\displaystyle\frac{B_{n-1}}{n-1}\right)=\mathrm{denom}\left(\displaystyle \frac{A}{(n-1) D}\right)=\frac{ \left(n-1\right)D}{\gcd(A,  n-1)}.
\]
On the other hand,
\[
\mathrm{denom}\left(\displaystyle B_{n-1}+\frac{1}{n}\right)=\mathrm{denom}\left(\displaystyle \frac{An+D}{D\cdot n}\right)=\frac{D\cdot n}{\gcd(An+D,  D\cdot n)}.
\]
Hence:
\[
\tilde{\chi}(n)=\frac{(n-1)\gcd(An+D, D\cdot n)}{n\gcd(A, n-1)}.
\]
It is evident that $\gcd(A, n-1)$ and $n$ are coprime. Therefore, the integrality condition for $\tilde{\chi}(n)$ is equivalent to the fact that $n \mid \gcd\left(An+D, D\cdot n\right)$. But $\gcd\left(An+D, D\cdot n\right)$ is divisible by $n$ if and only if $n\mid \left(An+D\right)$. And this, in turn, is equivalent to  $n\mid D$. 
In summary, we can say that $\tilde{\chi}(n)\in\mathbb{Z}$ if and only if $n\mid D$. 
\end{proof}


It should be noted how the expression for the function $\tilde{\chi}(n)$  can be simplified, if its argument is any Carmichael number.   We can write $D=\mathrm{denom}\left(B_{n-1}\right)=n\cdot m$, so that $\gcd\left(n, m\right)=1$. We get in succession:
\[
B_{n-1}+\frac{1}{n}=\frac{A+m}{nm}\;\; \Leftrightarrow \;\;  \mathrm{denom}\left(B_{n-1}+\frac{1}{n}\right)=\frac{n\cdot m}{\gcd(A+m, n)}\;\; \Leftrightarrow 
\]
\begin{equation}
\tilde{\chi}(n)=\frac{\left(n-1\right)\gcd(A+m, n)}{\gcd(A, n-1)}.
\label{6753867}
\end{equation}
\begin{lemma}
If for some Carmichael number $n$ we have  $n \mid \tilde{\chi}(n)$ then $n$ is a Giuga number.  
\end{lemma}
\begin{proof}
To begin with, we show that the condition $n\mid \left(A+m\right)$ is equivalent to the Giuga number condition. More exactly: 
\begin{equation}
n\mid \left(A+m\right) \Leftrightarrow n \equiv p\; \left(\mathrm{mod}\; p^2\right),\; \forall p\mid n.
\label{7647889}
\end{equation}

By the von Staudt-Clausen theorem, for a Carmichael number, the congruence 
\[
p\cdot B_{n-1}\equiv -1\; (\mathrm{mod}\; p),\;\; \forall p\mid n
\] 
holds.   The following logical chain of congruences is valid:
\[
\frac{p A}{nm}\equiv -1\; (\mathrm{mod}\; p)\; \Leftrightarrow A \equiv -\frac{n}{p}m\; (\mathrm{mod}\; p) \Leftrightarrow 
A+m \equiv m\left(1-\frac{n}{p}\right)\; (\mathrm{mod}\; p).
\]
Consequently, it further follows that 
\[
p\mid \left(A+m\right) \Leftrightarrow \frac{n}{p} \equiv 1 \; (\mathrm{mod}\; p) \Leftrightarrow n \equiv p \; (\mathrm{mod}\; p^2).
\]
Since these arguments hold for any divisor of $n$, and since $n$ is squarefree, it follows that it implies (\ref{7647889}). Thus, we have shown the equivalence of condition $n\mid \left(A+m\right)$ and the condition for a Giuga number (\ref{867434}). 

It follows, from (\ref{6753867}) that 
\[
n\mid \tilde{\chi}(n) \Leftrightarrow n\mid \gcd (A+m, n) \Leftrightarrow n\mid \left(A+m\right) \Leftrightarrow n \equiv p\; \left(\mathrm{mod}\; p^2\right),\; \forall p\mid n.
\]
Thus, the lemma is proved.
\end{proof}

Combining the lemmas stated above, we obtain the following result.
\begin{theorem}
The function $\chi(n)$ is integer-valued by $n$ if and only if $n$ is either a prime number or an odd Giuga number.
\end{theorem}



\section{Empirical observations and remarks}

It should be noted that the constancy condition $\chi(n)=\lambda$ yields a partition of the set of prime numbers into disjoint classes. 
Let us denote, say, by $P_{\lambda}$ the class of prime numbers for which the function $\chi(p)$ takes  the value $\lambda$. According to Lemma \ref{875643}, this number is even.
Let us consider the following case.  The sequence of even numbers beginning as $(2, 4, 6, 12, 16, 18, 36, 42,\ldots)$ is characterized by the fact that the first prime number in the sequence $P_{\lambda}$ begins with $\lambda+1$.  We suppose that this set is infinite.
In the following table, we present the first ten prime numbers belonging to the class $P_{\lambda}$.
\begin{center}
\begin{tabular}{|c|c|c|c|c|c|c|c|c|c|c|}
\hline
$P_2$ &3 & 11 & 23 & 47 & 59 & 71 & 83 & 107 & 131 & 167 \\
\hline
$P_4$ & 5 & 29 & 53 & 149 & 173 & 197   & 269 & 293 & 317  & 389 \\
\hline
$P_6$ & 7 & 31 & 67 & 79 & 103  & 139  & 151 & 223  & 283 & 367 \\
\hline
$P_{12}$ & 13 & 229  & 277 & 349  & 373 & 709 & 733 & 853 & 877 & 997 \\
\hline
$P_{16}$ & 17 & 113 & 593 & 977 & 1553 & 2129 & 2417 & 2609 & 2897 & 3089  \\
\hline
$P_{18}$   & 19 & 199  & 307 & 523 & 739 & 1063 & 1171 & 1279 & 1423 & 1531 \\
\hline
$P_{36}$ &  37   & 397 & 613 & 829 & 1117 & 1549 & 1693 & 2557 & 3637 & 3709 \\
\hline
$P_{40}$  &  41 & 281 & 521 & 761 & 1481  &  1721 & 2441  & 3881 & 5081 & 6521 \\
\hline
$P_{42}$  &  43 & 211 & 463  & 547 & 967  &  1051   & 1303  & 1723 & 2731 & 3067\\
\hline
 $P_{60}$ & 61  & 1021 &  1381 & 1741 &  2221 &  3181 &  3541 & 4021 & 4261  & 5821\\
\hline
$P_{72}$ & 73 &  1657 & 2089  & 2953 &  5113 &  5689 & 7417  & 8713 & 9433  & 10009 \\
\hline
$P_{96}$ & 97 & 4129  & 4513  & 5857 &  9697 & 10273  & 11617  & 12577 & 15073  & 15649 \\
\hline
$P_{100}$ & 101 & 701  & 1301  & 1901 & 3701  & 6101  & 6701 & 7901 & 10301  & 13901 \\
\hline
\end{tabular}
\end{center}
One sees, that there exist 13 such classes $P_{\lambda}$, where even number  $\lambda$ ranges from 2 to 100.


It looks like the condition $\chi(p)=p-1$ selects an integer sequence (excluding the first term, 2) known in the OEIS as \textrm{A337119}. What definition does the OEIS provide for this sequence? The definition states that these are all primes $p$ such that 
\[
b^{p-1} \equiv 1\; (\mathrm{mod}\; p-1)
\] 
for every $b$ coprime to $p-1$. 

%It is noteworthy that the sequence \textrm{A337119} is related to the sequence \textrm{A124240}.



Let us now consider $P_2$. It seems to form the sequence \textrm{A141123} (excluding the initial term, 2). If so, then these are precisely the primes representable by the quadratic form 
$F(x, y)=-x^2+2xy+2y^2$ of discriminant $D=12$.  
\begin{conjecture}
For any prime $p$: 
\[
\frac{p-1}{2}\;\; \mbox{divide}\;\; \mathrm{num}(B_{p-1})\; \Leftrightarrow\;  p\;\; \mbox{is}\;\; \mbox{representable}\;\; \mbox{by}\;\; -x^2+2xy+2y^2.
\] 
\end{conjecture}
We encourage the reader to prove or disprove this conjecture.

\section*{Acknowledgments}

This work was carried out within the framework of  the state assignment of the Ministry of Education
and Science of the Russian Federation on the project No.  126021217175-3.


\begin{thebibliography}{9}

\bibitem{Borwein1}
D. Borwein, J. M. Borwein, P. B. Borwein, and R. Girgensohn, 
Giuga's Conjecture on Primality,
\textit{ Amer. Math. Monthly} 
\textbf{103} (1996), 40-50.

\bibitem{Borwein}
J. Borwein, C. Maitland,  and M. Skerritt,   
Computation of an improved lower bound to Giuga's primality conjecture, 
\textit{Integers}
\textbf{13} (2013), Paper A67, 14 pp.

\bibitem{Clausen}
T. Clausen,  
Lehrsatz aus einer Abhandlung \"uber die Bernoullischen 
\textit{Zahlen. Astr. Nachr.} 
\textbf{17} (1840), 351-352.

\bibitem{Crandall}
R. Crandall, and C.  Pomerance,  
Prime Numbers: A Computational Perspective, 
Springer-Verlag, 2001.

\bibitem{Hardy}
G. H. Hardy, and  E. M. Wright,   
An Introduction to the Theory of Numbers, 6th ed. Oxford University Press, 2008.



\bibitem{Kellner1}
B. C. Kellner,  
The equivalence of Giuga's and Agoh's conjectures, 
\textit{arXiv  math/0409259} 
(2004).

%\bibitem{Kellner2}
%B. C. Kellner, and J. Sondow, 
%On Carmichael and polygonal numbers, Bernoulli polynomials, and sums of base-$p$ digits
%\textit{Integers} 
%\textbf{21} 
%(2021), 
%Paper A52, 21 pp.

\bibitem{Korselt}
A. Korselt,   
Probl\`eme chinois,  
\textit{L'interm\'ediaire des math\'ematiciens} 
\textbf{6} (1899)  142-143.

\bibitem{Staudt}
K. G. C. von Staudt,  
Beweis eines Lehrsatzes die Bernoulli'schen Zahlen betreffend, 
\textit{J. Reine Angew. Math.} 
\textbf{21} 
(1840), 
372-374.

\end{thebibliography}


\end{document}
